\chapter{Introduction}
\label{chap:intro}

\section{Context}
The job market is a complex ecosystem where both sides, job seekers and recruiters, try to find the best match for their needs. Unlike markets like Netflix, the job market has unique characteristics where only 1 job can be assigned to 1 person. The digitization should in theory make this process more efficient, enabling rapid and accurate matching. However, in practice, the job market suffers from this digital transformation due to its characteristics. While accessing information has become easier, the efficiency of matching decreased due to market congestion, information overload and systemic algorithmic biases. This thesis investigates how AI, more specifically \gls{RS}, can be used to improve matching quality in this two-side market by reducing congestion and improving recommendation relevance for both job seekers and recruiters.
\section{Motivation}
\label{sec:intro:motivation}

As more and more job search and recruitment processes are taking place online, the need for efficient and effective \gls{RS} becomes increasingly important. The need for a fair and efficient two-sided recommendation system forms a foundation for a more balanced job market. 
\vspace{1em}
In digital platfoms, visibility is key for success. Recommendation systems suffer from the problem of market congestion, where popular items (jobs) receive more attention compared to new or less popular items. 


\section{Problem Statement}
\label{sec:intro:problem}

Having set the stage for the importance of optimizing the labor market through improved recommendation systems, it is important to clearly define the specific technical challenges this thesis aims to address. The problem consists of 2 parts:

\begin{itemize}
	\item \textbf{For job seekers:} The current recommendation systems often fail to provide relevant job recommendations due to the overwhelming amount of information and the presence of popular jobs that overshadow less popular but potentially more suitable or probable opportunities. This leads to a suboptimal user experience and can result in job seekers missing out on relevant opportunities.
	\item \textbf{For recruiters:} Recruiters face the challenge of efficiently identifying suitable candidates from a large pool of applicants. The current systems may not effectively prioritize candidates based on relevance, leading to increased time and effort in the recruitment process and potentially overlooking qualified candidates.
\end{itemize}

Summarized, there's need for a \gls{RS} which provides relevant recommendations for both sides and mitigates the problem of market congestion, ensuring that both job seekers and recruiters can efficiently find suitable matches without being overwhelmed by irrelevant information or overshadowed by popular items.



\section{Research Objectives}
\label{sec:intro:objectives}

Outline the development of a two-way recommendation system to mitigate these imbalances and congestion.
